\section[Dynamic Windows]{Dynamic Windows \newmark}\label{sec:hmi-dynamic-windows}
Every equipment template opens a dynamic window: a panel on the right of the screen that shows the device's state, mode, commands, permissives, interlocks, alarms and history. Equipment of the same kind shares one dynamic window (Table~\ref{tab:hmi-dw-ids}), so a valve, motor or loop looks and works the same on every screen.

\begin{small}
\begin{longtable}{@{}P{3.6cm}P{12.6cm}@{}}
\caption{Dynamic Windows}\label{tab:hmi-dw-ids} \\
\toprule
\rowcolor{tblhead}\textbf{Window id} & \textbf{Used by} \\
\midrule
\endfirsthead
\multicolumn{2}{@{}l}{\small\textit{Table \thetable{} continued}} \\
\toprule
\rowcolor{tblhead}\textbf{Window id} & \textbf{Used by} \\
\midrule
\endhead
\texttt{dw\_ValveDiscrete} & On/off valves (\texttt{T\_DValve}) \\
\texttt{dw\_ValveAnalog} & Control valves (\texttt{T\_CValve}) \\
\texttt{dw\_Rotating} & Fixed-speed motors, pumps, fans, agitators and compressors (\texttt{T\_Motor}) \\
\texttt{dw\_VFD} & Variable-speed drives (\texttt{T\_VFD}) \\
\texttt{dw\_Heater} & Electric heaters (\texttt{T\_Heater}) \\
\texttt{dw\_HeatTransfer} & Heat exchangers (\texttt{T\_HeatEx}) \\
\texttt{dw\_ProcessValue} & Analog process values and indicators (\texttt{T\_AI}) \\
\texttt{dw\_Loop} & PID loops (\texttt{T\_PID}), including tuning \\
\texttt{dw\_Station} & Fill stations \\
\texttt{dw\_IOCard} & PLC I/O cards and racks (diagnostics) \\
\bottomrule
\end{longtable}
\end{small}

A template's \texttt{dockId} parameter overrides the default window chosen from its \texttt{tType} parameter, and \texttt{dwTagPath} overrides the bound tag. The \texttt{native} type opens the I/O card window; I/O card, rack and power-supply faceplates bind to the Modicon x80 device DDTs (decision 23).

I/O card, rack and power-supply faceplates bind directly to the Modicon x80 device DDTs from Control Expert, which are imported into Ignition as-is; there is no PCLib type for them (\texttt{T\_IOCard} and \texttt{T\_IORack} are removed), and the window id \texttt{dw\_IOCard} is kept.

\subsection[Size, Opening and Closing]{Size, Opening and Closing \newmark}\label{sec:hmi-dw-open}
\begin{itemize}[itemsep=1pt]
  \item \textbf{Size and position.} The dynamic window fills the full height between the header and the alarm bar, in the right-hand region of the shell: 400 $\times$ 860 at 1920 $\times$ 1080 (37.04\,vh wide, 79.63\,vh tall), and 800 $\times$ 1720 at 3840 $\times$ 2160 (Appendix~\ref{app:screen-design}).
  \item \textbf{Opening.} A single click or tap on the template symbol opens its dynamic window. The click target is the whole symbol bounding box, at least 44 design px in each direction.
  \item \textbf{Closing.} The close button returns the right-hand region to the screen's default view. Each screen view has a \texttt{defaultDock} view parameter that names it (normally the screen's trends), and \texttt{setDefaultDock} sets it on startup.
\end{itemize}

\subsection[Layout and Buttons]{Layout and Buttons \newmark}\label{sec:hmi-dw-layout}
Sections appear in this order; a section that does not apply to the device is left out:
\begin{enumerate}[itemsep=1pt]
  \item Header: tag, description and close button.
  \item State.
  \item Mode.
  \item Commands.
  \item Permissives.
  \item Interlocks.
  \item Alarms: active alarms are listed; inactive alarms are folded.
  \item History: an 8~h state strip and the most recent events.
  \item Maintenance and Configuration: collapsed by default and shown only to the roles in Table~\ref{tab:hmi-roles}. Loop tuning is done here, in \texttt{dw\_Loop}; there is no separate tuning popup.
\end{enumerate}
Buttons come in two sizes so that every dynamic window looks the same:
\begin{itemize}[itemsep=1pt]
  \item \textbf{Command tier} (mode AUTO /\allowbreak{} MAN /\allowbreak{} MAINT /\allowbreak{} OOS, OPEN /\allowbreak{} CLOSE, START /\allowbreak{} STOP, RESET, JOG, FWD /\allowbreak{} REV): one size, 72 $\times$ 30 design px, style \texttt{Button/Default}. The button for the current state is highlighted with \texttt{Button/Toggle/On}.
  \item \textbf{Secondary tier} (BYPASS, UNBYP, ACK, RESET LATCHED, RESET MAINT, SIMULATE, analog valve presets 0 /\allowbreak{} 25 /\allowbreak{} 50 /\allowbreak{} 75 /\allowbreak{} 100\,\%): one smaller size, style \texttt{Button/Small}.
  \item \texttt{Button/Danger} is used only for the confirm action in the bypass confirmation (Section~\ref{sec:hmi-bypass}).
\end{itemize}

\subsection[Touch Operation]{Touch Operation \newmark}\label{sec:hmi-touch}
In touch mode (template parameter \texttt{useButton = false}) a double tap on the symbol of a two-state device (\texttt{T\_DValve}, \texttt{T\_Motor}) sends the toggle command (open /\allowbreak{} close, start /\allowbreak{} stop). It is sent only when the device is in Manual and available, and is not interlocked, in local or faulted; otherwise the double tap does nothing. Each template has a \texttt{dblCmd} parameter (default on) to turn this off. Analog devices (control valves, drives, loops) take commands from the dynamic window only.

\subsection[Permissives, Interlocks and Bypass]{Permissives, Interlocks and Bypass \newmark}\label{sec:hmi-bypass}
The Permissives and Interlocks sections list each condition by name (\texttt{Cfg\_CondDesc} in \texttt{T\_Intlk}, Section~\ref{sec:pclib-t-intlk}) with its state, and mark the first-out, latched and bypassed conditions.
\begin{itemize}[itemsep=1pt]
  \item \textbf{Who.} Bypass and unbypass need the Supervisor role or higher. The BYPASS button is shown only for conditions configured as bypassable (\texttt{Cfg\_Bypassable}) and works only while the bypass authorisation \texttt{Inp\_BypassPerm} is on.
  \item \textbf{Confirmation.} BYPASS opens a confirmation (\texttt{Docks/Parts/BypassConfirm}). The user picks a duration from the presets 1~h, 2~h, 4~h and 8~h, limited to the group's maximum (\texttt{Cfg\_BypassMax}), and must enter a reason. The bypass is not applied without a reason.
  \item \textbf{Indication.} Each bypassed condition shows its own time left and an audit line: who bypassed it, their role, when, and the reason.
  \item \textbf{Logging.} Bypass and unbypass are written to the Ignition audit log (Section~\ref{sec:ignition-roles}).
  \item \textbf{Rule.} A bypass made this way from the HMI, on a condition configured as bypassable, is exempt from the second-person and approved MOC requirement in Section~\ref{sec:online-changes}. Safety-related signals are never configured as bypassable and cannot be bypassed in software. PLC forces and logic changes are not exempt.
\end{itemize}
\future{Named permissive and interlock groups are to be embedded in the device types: \texttt{Perm : T\_Intlk} and \texttt{Intlk : T\_Intlk} in \texttt{T\_DValve}, \texttt{T\_Motor} and \texttt{T\_VFD}, \texttt{Intlk : T\_Intlk} in \texttt{T\_CValve} and \texttt{T\_PID}, and \texttt{Prot : T\_Intlk} in equipment modules. The existing \texttt{Inp\_IntlkOK} and \texttt{Inp\_PermOK} are kept and ANDed with the group. \texttt{T\_Intlk} gains \texttt{Sts\_BypLeft : ARRAY[0..15] OF TIME} so the time left can be shown per condition; until then only the time left on the oldest bypass (\texttt{Sts\_BypassLeft}) is available.}

\subsection[Trends]{Trends \newmark}\label{sec:hmi-trends}
Trends in a dynamic window show the last 30~minutes by default. Every graph has an expand button. Expanding a graph keeps the dynamic window at 400 $\times$ 860 design px and the right-hand region never changes width: the graph fills the window below its header (about 400 $\times$ 760 design px) with range buttons 5~m, 30~m, 1~h, 8~h and 24~h, pause /\allowbreak{} live and a collapse button. The other sections of the window are hidden until the graph is collapsed.
